\documentclass{article}
\usepackage{iclr2024_conference,times}

\input{math_commands.tex}

\usepackage{booktabs}
\usepackage{hyperref}
\usepackage{url}
\usepackage{graphicx}
\usepackage{multirow}
\usepackage{xcolor}
\usepackage{enumitem}
\usepackage{float}
\renewcommand{\ttdefault}{cmtt}

\iclrfinalcopy

\title{Study Coach: R3. Modification --- Single-Task Feedback-Only Prompt}

\author{
  Dipti Aswath \hspace{2em} Buddhika Makehelwala \hspace{2em} Jiashan Cui \hspace{2em} Khubi Shah \\
  \texttt{\{daswath, bmakehelwala, jiashanc, khubis\}@andrew.cmu.edu}
  }

\date{}

\begin{document}
\maketitle

\paragraph{Jiashan Contribute.}
Implemented the feedback-only prompt modification, decoupling coaching feedback generation from verdict classification. Ran the multimodal C4 evaluation using the feedback-only prompt on both Qwen3-VL-8B-Instruct and Qwen3-VL-32B-Instruct ($N{=}50$ each), collecting LLM-as-Judge and automated metric (F1, ROUGE-L, BLEU) results.

\section{Modification: Decoupling Feedback Generation from Classification}

\subsection{Motivation}

A central finding from Report 2 was a verdict--feedback dissociation (H2): the multimodal condition (C4)
produced the best coaching feedback but the \emph{worst} verdict accuracy across both model scales.
Specifically, text-only C1 achieved 56\% verdict accuracy at both 8B and 32B, while C4 achieved only
48\% (8B) and 54\% (32B). Meanwhile, C4 led all conditions on feedback quality: 6\% LLM-judge match and
80\% human match at 8B, and 20\% LLM-judge match at 32B.

A classmate observation during report feedback directly motivated our modification:
\begin{quote}
\emph{``Your result that shows text-only wins classification but multimodal wins feedback quality is
really useful. Have you thought about doing it in 2 stages? Like text only first and then multimodal
grounding to get the best of both?''}
\end{quote}

This observation pointed to a structural problem: our H1/H2 prompt asked the model to simultaneously
produce three outputs---a verdict (Correct / Partially Correct / Incorrect), an error category
(Factual / Conceptual / Omission), and free-text coaching feedback---in a single pass with a single
prompt. Under this multi-task framing, the classification objective and the explanation objective
compete for the model's capacity, and the verdict structure in the prompt likely anchors the model
toward classification-style reasoning even when generating feedback.

The classmate's two-stage intuition suggested treating grading and coaching as \emph{separate concerns}.
As a first step toward validating this architecture, we simplified the prompt to a single objective:
\textbf{generate coaching feedback only}, with no reference answer and no verdict or error category
requested. This isolates the multimodal grounding capability for feedback generation, and tests
whether removing the classification scaffolding allows the model to produce higher-quality explanations.

\subsection{Experimental Setup}

We applied the feedback-only prompt to the multimodal condition (C4: question + figure image + caption)
on both the baseline and competitive models:

\begin{itemize}[itemsep=4pt]
    \item \textbf{Baseline}: Qwen3-VL-8B-Instruct, C4 multimodal, $N{=}50$
    \item \textbf{Competitive}: Qwen3-VL-32B-Instruct, C4 multimodal, $N{=}50$
\end{itemize}

Both runs use the same evaluation indices and are judged by Claude Sonnet 4 (LLM-as-Judge), enabling
direct comparison with the H1/H2 multi-task baseline results reported in Table~\ref{tab:h1h2_r2}.
The prompt template is reduced from:

\begin{quote}
\texttt{Verdict = \ldots} \\
\texttt{Error Category = \ldots} \\
\texttt{Feedback = \ldots}
\end{quote}

to simply:

\begin{quote}
\texttt{Feedback = \textless concise study coach explanation, second-person tone\textgreater}
\end{quote}

No ground-truth answer is provided in either condition, consistent with our no-reference-answer
experimental design (providing the answer was found in R2 to reduce modality differences and lower
alignment with human judgment).

\subsection{Results}

Table~\ref{tab:modification} reports LLM-as-Judge results (N=50) and automated metrics (token-level F1,
ROUGE-L, BLEU) for the multi-task H1/H2 prompt versus the feedback-only prompt, at both model scales.

\begin{table}[H]
\centering
\small
\caption{Feedback quality: multi-task H1/H2 prompt vs.\ feedback-only prompt, C4 multimodal,
$N{=}50$. LLM-as-Judge labels: Match / Partial / Unmatched. Soft Match = Match + Partial.}
\label{tab:modification}
\begin{tabular}{@{}llrrrrrrrr@{}}
\toprule
Model & Prompt & Match & Partial & Unmatch & Match\% & Soft\% & F1 & ROUGE-L & BLEU \\
\midrule
\multirow{2}{*}{Qwen3-VL-8B}
 & Multi-task (H1/H2) & 3  & 25 & 22 & 6.0\%  & 56.0\% & 0.302 & 0.197 & 6.16 \\
 & Feedback-only      & 4  & 29 & 17 & 8.0\%  & \textbf{66.0\%} & \textbf{0.377} & \textbf{0.242} & \textbf{7.30} \\
\midrule
\multirow{2}{*}{Qwen3-VL-32B}
 & Multi-task (H1/H2) & 10 & 19 & 21 & \textbf{20.0\%} & 58.0\% & 0.298 & 0.202 & 7.44 \\
 & Feedback-only      & 9  & 29 & 12 & 18.0\% & \textbf{76.0\%} & \textbf{0.404} & \textbf{0.257} & \textbf{10.45} \\
\bottomrule
\end{tabular}
\end{table}

\paragraph{8B results.}
The feedback-only prompt improves soft match from 56\% to 66\% ($+$10pp) and reduces unmatched cases
from 22 to 17 ($-$23\%). All automated metrics improve substantially: F1 $+$25\%, ROUGE-L $+$23\%,
BLEU $+$18\%. Partial matches increase (25 $\rightarrow$ 29), reflecting that more responses are
directionally correct even if not yet fully complete. Strict match improves modestly (6\%
$\rightarrow$ 8\%).

\paragraph{32B results.}
The improvement is larger at 32B scale. Soft match rises from 58\% to 76\% ($+$18pp) and unmatched
cases drop from 21 to 12 ($-$43\%). Automated metrics again improve across the board: F1 $+$36\%,
ROUGE-L $+$27\%, BLEU $+$40\%. Partial matches increase substantially (19 $\rightarrow$ 29), and
strict match is essentially unchanged (20\% $\rightarrow$ 18\%), suggesting the 32B model already
achieves strong binary alignment under the multi-task prompt, and the gain from task simplification
lies primarily in improving responses that were previously ``directionally correct but incomplete.''

\paragraph{Scale comparison.}
Table~\ref{tab:scale} compares the two models under the feedback-only prompt directly.

\begin{table}[H]
\centering
\small
\caption{Feedback-only prompt: 8B vs.\ 32B comparison, C4 multimodal, $N{=}50$.}
\label{tab:scale}
\begin{tabular}{@{}lrrrrrrrr@{}}
\toprule
Model & Match & Partial & Unmatch & Match\% & Soft\% & F1 & ROUGE-L & BLEU \\
\midrule
Qwen3-VL-8B  & 4 & 29 & 17 & 4.0\%  & 66.0\% & 0.377 & 0.242 & 7.30 \\
Qwen3-VL-32B & 9 & 29 & 12 & \textbf{18.0\%} & \textbf{76.0\%} & \textbf{0.404} & \textbf{0.257} & \textbf{10.45} \\
\bottomrule
\end{tabular}
\end{table}

Scaling from 8B to 32B with the feedback-only prompt yields consistent gains across all metrics. Strict
match improves from 4\% to 18\% ($+$14pp) and soft match from 66\% to 76\% ($+$10pp), with unmatched
cases dropping from 17 to 12. The partial match count is identical (29 each), suggesting the ceiling on
``directionally correct'' responses is model-scale-independent --- the 32B advantage lies in converting
partial responses into full matches rather than reducing errors from the bottom.

\subsection{Analysis}

\paragraph{Better prompting matters even at larger scale.}
The 32B multi-task baseline (58\% soft match) is actually \emph{worse} than the 8B feedback-only
result (66\% soft match), despite a 4$\times$ parameter increase. The 32B feedback-only configuration
(76\% soft match, F1 = 0.404) is the best result across all configurations and both scales. This
demonstrates that task simplification is not just a workaround for small-model limitations --- it
provides a consistent and additive benefit regardless of model capacity.

\paragraph{The verdict--feedback dissociation persists and motivates separation.}
Our R2 finding that classification and explanation objectives are inversely correlated (H2) is further
supported here: removing the classification requirement does not hurt strict match meaningfully (within
2pp at both scales) and substantially improves all measures of explanation quality. This is consistent
with the hypothesis that the multi-task prompt creates a classification-anchoring effect that
constrains how the model generates feedback.

\paragraph{Connection to the two-stage architecture proposal.}
These results provide empirical motivation for the condition-routing architecture proposed in R2: use
text-only (C1) for grading and a feedback-focused multimodal prompt for coaching. The feedback-only
modification is the coaching branch of this two-stage system, validated at both 8B and 32B scale.
The next step would be to implement the classification stage (text-only or caption-only) and combine
verdicts with feedback-only outputs into a unified StudyCoach response, testing whether the two-stage
system closes the verdict--feedback dissociation gap without sacrificing explanation quality.

\end{document}
